\chapter{绪论}\label{chap:introduction}

\section{财富水平塑造经济决策}\label{chap1:utility_theory}
财富水平如何影响经济决策，是神经经济学关注的核心问题之一。1738年，Daniel Bernoulli首次尝试以数学公式量化财富水平与经济决策之间的关系\cite{bernoulli1954exposition}。为解释“圣彼得堡悖论”，他提出人们最大化的并非货币本身，而是货币所承载的主观价值——即效用（附录一）。在对数形式效用函数$u(w)=\log(w)$下（图~\ref{fig0_1:marginal_utility}），同等货币增量$\Delta w$带来的收益$\Delta u$并非恒定，而是取决于初始财富水平：低财富区间的效用曲线陡峭，边际效用递减显著，个体的风险规避倾向更强；高财富者的效用曲线相对平缓，行为更接近理性模型预测的风险中性。这一框架后被完善并公理化为VNM效用理论\cite{vonneumann1953theory}（附录二）。该理论从一组关于偏好的理性公理出发，证明满足公理的决策者必然存在一个效用函数，使其行为等价于最大化期望效用，从而将效用从经验性描述提升为可检验的理论命题。在此基础上，前景理论\cite{kahneman1979prospect}进一步引入动态参考点与损失厌恶等机制，以处理效用框架的若干系统性偏离——例如人们有时宁愿接受绝对收益较低、但在社会比较中相对占优的结果\cite{solnick1998more,hill2010risk}。这类现象表明，在涉及财富的决策中，个体关心的不仅是绝对财富水平，更是在参照群体中的相对位置。以效用函数连接财富与决策行为，已成为经济决策研究的标准范式。

\begin{figure}[H]
    \centering
    \includegraphics[width=0.75\textwidth]{Figure0_1_Marginal_Utility.png}
    \bicaption{\enspace 凹效用函数与边际效用递减}{\enspace Concave utility function and diminishing marginal utility}
    \label{fig0_1:marginal_utility}
\end{figure}
\fignotetext{横轴为财富水平（$w$），纵轴为效用（$u$），即决策者持有财富$w$时的主观满意度。在凹函数设定下，效用函数的斜率单调递减，即边际效用随财富增加而下降。这意味着相同的财富增量$\Delta w$所带来的效用增量$\Delta u$并不恒定：在低财富区间（图左侧，深色）中的$\Delta u$显著大于高财富区间（图右侧，浅色）中的对应增量。}{The horizontal axis represents wealth level ($w$) and the vertical axis represents utility ($u$), interpreted as the decision-maker's subjective satisfaction from holding wealth $w$. Under a concave utility function, the slope of utility function decreases monotonically as wealth accumulates, implying that marginal utility diminishes as wealth increases. Consequently, an identical wealth increment $\Delta w$ yields different utility gain $\Delta u$ at different wealth levels: the increment in the low-wealth region (left, darker gray) substantially exceeds that in the high-wealth region (right, lighter gray).}

大量人类行为研究为上述理论框架提供了实证支持。例如，瑞典彩票中奖者及其配偶在财富骤增后工作时间显著下降\cite{cesarini2017effect}。这与边际效用递减的预测一致：财富的增加降低了劳动收入的边际效用，工作意愿随之减弱。对76国逾8万人的调查表明，富人表现出更高的利他倾向\cite{vanags2025greater}。这也可以在效用框架下得到解释：当财富的边际效用随累积而下降时，亲社会行为所附带的非财富收益——如声誉与社会关系——的相对价值上升，利他行为因此更具吸引力。值得注意的是，财富水平并非只影响成人的决策偏好，4--5岁儿童的风险偏好已受其社会经济地位的影响\cite{harvey2022developmental}。

这一问题在现实中有着深刻的社会意义。效用理论将穷人与富人的决策差异归结为在效用曲线上所处区间的不同——这是一种理性适应。然而，一些研究指出，长期资源匮乏可能导致更深层的认知劣势：它不仅改变决策者在效用函数上的位置，还可能改变函数本身的形状，甚至因非理性而使效用框架完全失效。在Mani等人的经典研究中，无论是想象的还是真实的财务困境，都会损害个体的认知能力\cite{mani2013poverty}。Kaur等人则发现，通过提前发放部分薪资缓解工人的资金短缺，能有效释放此前被财务焦虑占据的认知资源，在显著提升计件工人生产效率的同时，也大幅降低了工人的操作错误率\cite{kaur2025financial}。Haushofer与Fehr整合相关研究后提出，贫困的长期压力会系统性地消耗认知资源与执行控制力，使个体难以执行长远最优决策，从而陷入贫困与认知受损相互加剧的恶性循环\cite{haushofer2014psychology}。这一负反馈模型为贫困的持久性提供了整合性解释，但其因果链条的稳健性仍存争议。大规模的随机对照试验发现财富水平的波动仅影响金钱相关决策，而认知能力不随短期财务波动而改变\cite{carvalho2016poverty}，这直接挑战了Mani等人的结论。Haushofer在最新综述中也强调，匮乏感影响经济决策的实证证据整体偏弱，具体机制仍需更严格的受控实验加以澄清\cite{haushofer2023psychology}。从神经机制层面理解财富如何影响经济决策，或可为解决上述争论提供更为直接的证据，也为探索如何提升贫困人群的决策质量提供新的视角。

\section{财富认知的跨物种基础}\label{chap1:animal_BHV}
效用理论在人类研究中已获得广泛支持，但人类研究存在固有局限：我们既不能在健康被试中对特定脑区实施侵入性操纵以探究因果关系，也难以将先天生物因素与后天文化习得完全剥离。动物研究在这些问题上具有不可替代的优势。一方面，严格受控的实验室条件使得单细胞记录、脑区失活、电刺激、光遗传等技术手段成为可能，为解析决策行为的神经机制创造了条件。另一方面，跨物种比较使我们得以追溯经济决策行为的演化根源——哪些计算原理是保守且先天的，哪些是人类在语言、文化与社会制度环境中衍生的特有产物。回答这些问题能帮助建立有生物学基础的经济决策理论。

尚无证据表明动物拥有与人类的货币财富系统同构的符号体系。但货币在本质上只是符号性财富的一种形式，其所承载的核心功能并不为货币所独有。探索“财富效应”的跨物种基础，可以考察动物能否理解其他形式的符号性财富的核心功能，如延迟支付与交换媒介\cite{addessi2020roots,beran2021primate}。在延迟支付层面，Pietras等人证实，经过系统性强化训练后，原本中性的刺激（如LED灯）便足以构成独立的奖惩机制，稳定驱动鸽子的行为\cite{pietras2005punishment}。这说明代币已被赋予独立的价值，成为有效的条件强化物。Seo等人在猕猴博弈任务中进一步发现，猕猴能根据代币的得失调节决策权重——获得代币显著强化当前选择，失去代币则产生抑制效应\cite{seo2009behavioral}。Sousa等人进一步发现黑猩猩在代币奖励任务中自发表现出“储蓄”倾向，其代币积累量远超单次兑换所需的最低数量\cite{sousa2001use}。以上实验共同表明，动物能够理解代币作为延迟支付凭证所承载的跨时间价值，并将其作为有意义的次级强化物加以使用。在交换媒介方面，Chen等人发现卷尾猴能根据商品价格波动灵活调整消费策略：他们会在市场波动时转向购买降价商品\cite{chen2006basic}。这表明它们能够将代币理解为一般等价物并与多种商品进行价值比较。但上述功能证据也存在一些局限，例如动物的储蓄行为能否延伸至更长时间尺度（如跨天）？代币交易能否被自然地扩展到猕猴生活环境中的所有资源交换中？虽然巴厘岛野生长尾猕猴表现出“以物换食”行为\cite{leca2021acquisition}，但这究竟是灵长类普遍具备的认知潜力，还是特定生态环境催生的文化特例，尚无定论。总体而言，虽然动物代币系统已能有效模拟货币的核心功能，但是与人类财富体系之间仍存在差距。

在理解"财富"的基础上，一系列研究进一步证实，累积的符号性财富能够影响恒河猴的经济决策。Yang等人发现猕猴在代币博弈任务中的决策具有显著的财富水平依赖性：在获益情景中随着代币资产累积，额外代币的边际效用相对下降，风险规避倾向随之增强（图~\ref{fig0_2:yang2022}）\cite{yang2022primate}。Ferro等人也发现类似的结论：当积累代币数趋近奖励阈值时，猕猴的经济策略趋于审慎——风险厌恶倾向随之增强\cite{ferro2026accumulation}。综上，非人灵长类在符号性财富情境下展现的决策模式与效用理论的预测高度一致。这不仅提示"财富效应"具有演化基础，也为进一步探究财富调节决策的神经机制提供了行为学依据。

\begin{figure}[H]
    \centering
    \includegraphics[width=\textwidth]{Figure0_2_Yang2022_Paradigm}
    \bicaption{\enspace 代币财富调节猕猴风险决策偏好}{\enspace Token wealth modulates risky preferences in macaques}
    \label{fig0_2:yang2022}
\end{figure}
\fignotetext{\textbf{(A) 代币赌博任务。}试次开始时，注视点周围出现实心圆点（代币），圆点数量表示当前持有的代币数。猕猴注视中心点0.5--1秒后，屏幕上同时呈现一个确定选项和一个赌博选项。猕猴通过眼跳注视所选目标以作出选择；之后代币移至该目标周围，未选目标消失。猕猴需继续注视所选目标450--550毫秒，随后目标外观发生变化以提示选择结果：赌博选项由双色方块变为单色方块，颜色表示代币的得失量；确定选项由方形变为圆形，作为视觉变化的对照。再经过500毫秒代币状态更新：若持有代币数达到6个或以上，则给予600微升饮用水奖励，超出6个的代币累积至下一试次；若代币数不足6个，则不给予奖励。\\\textbf{(B) 获益情境与损失情境下的选项集合。}所有选项均用色块表示，颜色代表代币的得失量（范围为$-3$至$+3$），色块面积比例代表该结果发生的概率。确定选项为单色方块，表示确定的代币得失。赌博选项为双色方块，包含两种可能结果。本任务共使用6种赌博选项：获益情景下可能得到3个或0个代币，三个选项的概率组合分别为$[0.1, 0.9]$、$[0.5, 0.5]$和$[0.75, 0.25]$；损失情景下可能失去0个或3个代币，概率组合与前述相同。每种情境中，3种赌博选项分别与4个确定选项配对，产生12种组合，总计24种选择组合。\\\textbf{(C) 选择赌博选项的概率P(Gamble)随起始代币数的变化。}在获益情境下（绿色），P(Gamble)随起始代币数增加而显著下降（猕猴G：$\beta = -0.044$, $p < 10^{-4}$；猕猴O：$\beta = -0.035$, $p < 10^{-4}$），表明代币财富积累越多，额外代币的边际效用递减，猕猴对进一步获益的风险规避越强。损失情境下（红色），P(Gamble)随起始代币数增加而上升，表现为风险寻求，符合前景理论的预测。图中数据以均值$\pm$标准误表示，黑色统计标记为各代币水平下的配对比较结果。改编自Yang et al. (2022)\cite{yang2022primate}。}{\textbf{(A) Token-based gambling task.} At the start of each trial, filled dots (tokens) appeared around the fixation point, with the number of dots indicating the currently held tokens. After the monkey maintained fixation for 0.5--1 s, a sure option and a gamble option were presented simultaneously on the screen. The monkey made its choice by shifting gaze to the desired target; the token cue then moved to surround the chosen target while the unchosen target disappeared. The monkey was required to keep fixating the chosen target for 450--550 ms, after which the target's appearance changed to reveal the outcome: a gamble option changed from a two-colored square to a single-colored square, with the color indicating the token gain or loss; a sure option changed from a square to a circle, serving as a control for the visual change. After another 500 ms, the token state was updated: if the token count reached six or more, a 600 $\mu L$ water reward was delivered and any tokens exceeding six were carried over to the next trial; if the count was below six, no reward was given.\\\textbf{(B) Option sets in the gain and loss contexts.} All options were represented by colored squares. Color denoted the magnitude of token gain or loss (ranging from $-3$ to $+3$), and the proportion of the colored area indicated the probability of that outcome. Sure options were single-colored squares representing a certain token gain or loss. Gamble options were two-colored squares involving two possible outcomes. Six gamble options were used: in the gain context, one could gain either 3 or 0 tokens, with the three options having probability combinations of $[0.1, 0.9]$, $[0.5, 0.5]$, and $[0.75, 0.25]$; in the loss context, one could lose either 0 or 3 tokens, with the same set of probability combinations. In each context, the three gamble options were paired with four sure options, yielding 12 combinations per context and 24 combinations in total.\\\textbf{(C) Probability of choosing the gamble option, P(Gamble), as a function of start token number.} In the gain context (green), P(Gamble) decreased significantly as the start token number increased (Monkey G: $\beta = -0.044$, $p < 10^{-4}$; Monkey O: $\beta = -0.035$, $p < 10^{-4}$), indicating that greater token wealth was associated with stronger risk aversion for further gains, reflecting the diminishing marginal utility of additional tokens. In the loss context (red), P(Gamble) increased with the start token number, showing risk-seeking behavior, consistent with the predictions of prospect theory. Data are presented as mean $\pm$ SEM. Black statistical markers indicate pairwise comparisons at each token level. Adapted from Yang et al. (2022)\cite{yang2022primate}.}

\section{经济决策与财富的神经机制}\label{chap1:neural_mechanism}
要理解财富如何调节价值决策，首先需要厘清价值决策本身的神经基础。本节依次介绍三个与之密切相关的脑区——眶额叶皮层（orbitofrontal cortex，OFC，\ref{chap1:OFC}节）、背外侧前额叶皮层（dorsolateral prefrontal cortex，DLPFC，\ref{chap1:DLPFC}节）与前扣带皮层（anterior cingulate cortex，ACC，\ref{chap1:ACC}节），随后聚焦于财富的神经表征及其对价值编码的调控作用（\ref{chap1:wealth_neural}节）。

\subsection{OFC：主观价值}\label{chap1:OFC}
OFC处于情感与认知的交汇点，在预测行为后果和指导最优决策中发挥着核心作用\cite{rudebeck2018orbitofrontal}。这一功能与其解剖连接模式高度吻合：OFC广泛汇聚来自嗅觉、味觉、视觉、躯体感觉及内脏感觉等多种通道的高级加工输入，并与杏仁核（amygdala）、海马体（hippocampus）等边缘结构保持紧密的双向联系\cite{cavada2000anatomical,passingham2012neurobiology}。这一连接模式为OFC将外部刺激与内部动机状态整合编码为统一的价值信号提供了解剖学基础。

OFC与腹内侧前额叶皮层（ventromedial prefrontal cortex，vmPFC）被认为是价值网络的核心节点，以跨类别的“神经通用货币”（neural common currency）形式整合不同来源的价值信号\cite{levy2012root}。来自人类功能磁共振成像（functional magnetic resonance imaging，fMRI）与猕猴电生理的证据共同支持了这一观点。例如，Plassmann等人发现，当同一种葡萄酒被赋予更高的标价时，被试报告的主观愉悦感显著增强，同时，其内侧眶额叶皮层（medial orbitofrontal cortex，mOFC）的神经活动也相应上升\cite{plassmann2008marketing}；这表明mOFC所编码的是物品的主观价值，而非其客观属性。此外，通过选择性饱食使特定奖励贬值后，人类和猕猴的OFC对该奖励的价值表征均显著减弱\cite{gottfried2003encoding,pastorbernier2021satiety}。这说明OFC的价值编码并非简单反映感觉输入，而是会依据个体当前的内部状态动态更新。在单神经元水平上，Padoa-Schioppa等人的系列实验首次系统揭示了OFC对主观价值的抽象编码。在其实验中，猕猴需要在不同口味和数量的果汁之间做出选择，这种决策被称为“基于物品的决策”，以区别于“基于动作的决策”。结果发现，无论选项呈现的位置或猕猴最终的眼动方向如何，OFC神经元都能稳定地编码选项价值（offer value）、被选价值（chosen value）以及被选类型（chosen juice）等抽象的“物品域”信号（图~\ref{fig0_3:padoa2006}）\cite{padoaschioppa2006neurons, padoaschioppa2011neurobiology, padoaschioppa2017orbitofrontal}。

\begin{figure}[H]
    \centering
    \includegraphics[width=0.85\textwidth]{Figure0_3_Padoa2006}
    \bicaption{\enspace 眶额叶皮层神经元编码主观价值}{\enspace OFC neurons encode subjective value}
    \label{fig0_3:padoa2006}
\end{figure}
\fignotetext{\textbf{(A) 实验范式。}猕猴注视屏幕中央1.5秒后，屏幕两侧各出现一组彩色方块，方块颜色代表果汁种类，数量代表果汁滴数（0--10滴）。在1--2秒的随机延迟之后，两个眼动目标出现，猕猴用眼动作出选择。继续注视所选目标0.75秒后获得对应的果汁。不同试次中，选项的左右位置进行平衡。\\\textbf{(B) 四个示例OFC神经元。}每一列为一个神经元，上图是该神经元记录当天的行为选择模式，下图是其调谐曲线。在上图中，散点及S形拟合曲线展示了不同选项组合下猕猴选择果汁B的比例，曲线拟合得出相对价值（如1A=3.3B表示1滴果汁A与3.3滴果汁B主观价值相等，此时猕猴选择两者的比例相同）；下图曲线所用的不同颜色代表不同分析时间窗口（蓝色为选项呈现之后，黑色为奖励发放之后，浅灰色为选项出现前的基线活动）；误差线表示标准误。这四个神经元分别编码：被选价值、果汁A的选项价值、果汁B的选项价值，以及被选果汁类型。改编自Padoa-Schioppa \& Assad, 2006\cite{padoaschioppa2006neurons}。}{\textbf{(A) Task paradigm.} After the monkey fixated the center of the screen for 1.5 s, two sets of colored squares appeared on either side of the fixation point. The color of the squares indicated the juice type, and the number of squares indicated the juice amount (0--10 drops). Following a random delay of 1--2 s, two saccade targets appeared, and the monkey indicated its choice with an eye movement. The monkey then maintained fixation on the chosen target for 0.75 s, after which the corresponding juice was delivered. The left/right positions of the offers were counterbalanced across trials.\\\textbf{(B) Four example OFC neurons.} Each column corresponds to one neuron, with the top panel showing the behavioral choice pattern from the recording day and the bottom panel showing the neuron's tuning curve. In the top panels, the sigmoid fit of the percentage of trials in which juice B was chosen across offer types yields a measure of relative value (e.g., 1A=3.3B means that 1 drop of juice A and 3.3 drops of juice B have equal subjective value, and the monkey chose the two options with equal probability). In the bottom panels, different colors indicate different analysis time windows (blue: post-offer; black: post-juice; light gray: pre-offer baseline). Error bars represent s.e.m. The four neurons encode, respectively: chosen value, the offer value of juice A, the offer value of juice B, and the chosen juice type. Adapted from Padoa-Schioppa \& Assad, 2006\cite{padoaschioppa2006neurons}.}

在单神经元分析的基础上，群体方法提供了新的视角。例如，McGinty等人发现，群体信号在试次间的波动能够预测动物的最终选择，而这一关系在单神经元水平上并不存在，提示价值比较依赖于群体编码而非单个神经元的稳定调谐\cite{mcginty2023behavioral}。类似地，Ferrari-Toniolo等人则发现，单神经元层面的分析可能会掩盖OFC编码的真实结构。单个OFC神经元对效用函数和概率权重函数的表征高度异质，很多神经元编码的风险态度与动物的实际行为并不一致。然而，当将这些异质性信号整合为群体编码时，OFC神经元群体却能形成与主观经济行为高度吻合的价值表征\cite{ferraritoniolo2023reliable}。在价值编码的时间动态上，Rich等人发现在决策过程中OFC的神经状态并非同时稳定地表征所有选项，而是在被选选项与未被选选项的价值表征之间动态交替，且这种状态切换的模式与决策行为密切相关：当被选选项的表征占主导时，决策速度更快；而当两种选项的表征强度相近时，被试的决策时间也更长。这表明，价值比较可能并非通过各选项价值的并行静态编码实现，而是依赖于OFC群体网络状态在时间上的动态切换\cite{rich2016decoding}。总体而言，OFC对主观价值的编码并非由单个神经元所决定，而是依赖于群体水平上的动态整合与分布式表征。

上述相关性研究表明OFC活动与主观价值或行为选择之间存在关联，但无法说明因果方向——这需要操纵性实验的检验。Ballesta等人利用电刺激直接扰动OFC活动，从两个层面揭示了因果作用：中等强度电流（约50微安）改变价值表征本身，可能借助OFC神经元的范围适应特性使得猕猴偏向价值范围更大的选项；而极低电流（5--15微安）在不改变价值表征的前提下干扰价值比较过程，从而在不改变偏好的情况下增加了选择变异性\cite{ballesta2020values, ballesta2022orbitofrontal}。综上，OFC是价值表征与价值比较的核心节点，对经济决策有因果性贡献。

\subsection{DLPFC：视觉空间框架与认知控制}\label{chap1:DLPFC}
DLPFC汇聚来自后顶叶皮层（posterior parietal cortex, PPC）、颞上沟（superior temporal sulcus, STS）及嗅周皮层（perirhinal cortex）的高级感觉输入，与OFC及海马系统保持紧密联系，并直接投射至前运动皮层（premotor cortex）与前辅助运动区（pre-supplementary motor area, pre-SMA）。这一连接模式使其成为整合多源信息并生成下行控制信号的枢纽\cite{passingham2012neurobiology}。

DLPFC的一项核心功能是在视觉空间框架中形成决策。Funahashi等人的奠基性研究首次发现，DLPFC神经元在延迟期内能够选择性地维持视觉线索的空间位置信息，形成方位特异性的“记忆野”（memory field，Funahashi等人\cite{funahashi1989mnemonic}）。在此基础上，Takeda等人采用旋转延迟眼动（rotatory oculomotor delayed-response, R-ODR）任务，将感觉信号与运动信号解离。群体向量分析表明，DLPFC的方向信息表征从最初表征视觉线索的方向逐渐旋转至表征眼动目标的方向\cite{takeda2004population}。这说明DLPFC并非简单复制感觉输入或预先激活运动指令，而是在延迟期内动态计算，逐步将感觉表征转化为运动规划。值得注意的是，DLPFC将信息锚定于空间坐标的特性，并不依赖于空间信号的输入。Lin等人在一项概率推理任务中，以完全抽象的价值信号（形状序列所携带的奖励概率）作为输入，发现DLPFC神经元在单条证据呈现时同时编码价值与动作信息，并能够跨时间累积动作域中的证据以形成最终判断；与此形成对比的是，眶额叶皮层仅编码单条证据的价值信息，且缺乏累积过程（图~\ref{fig0_4:lin2020}）\cite{lin2020evidence}。Cai等人在经济决策框架中得到了类似结论：在选项尚未呈现时，DLPFC在“物品域”（good-based domain）编码决策相关变量；在选项位置揭示后，则切换至“动作域”（action-based domain）编码，开始为即将发生的眼动做准备\cite{cai2014contributions}。上述研究共同表明，无论输入是空间信号还是抽象的价值信号，DLPFC均以具体的空间位置——如眼动方向——为决策计算的终点，其核心运算始终收敛于视觉空间坐标。

\begin{figure}[H]
    \centering
    \includegraphics[width=0.8\textwidth]{Figure0_4_Lin2020}
    \bicaption{\enspace OFC与DLPFC对概率证据的差异化编码}{\enspace Differential encoding of probabilistic evidence by OFC and DLPFC neurons}
    \label{fig0_4:lin2020}
\end{figure}
\fignotetext{\textbf{(A) 概率推理任务。}猴子注视中央注视点时，屏幕上依次呈现四个形状。每个试次中，四个形状从10个形状的集合中有放回地随机抽取，每个形状被赋予一个独特的权重值（见插图）。正权重意味着红色目标获得奖励的概率更高，负权重则意味着绿色目标获得奖励的概率更高。四个形状的权重之和决定了该试次中红、绿目标的相对奖励概率。注视点消失后，猴子通过眼跳选择红色或绿色目标，并依概率获得奖励。时间轴上红、黄、绿、蓝四色阴影标示四个形状依次呈现的时段。\\\textbf{(B) OFC与DLPFC示例神经元。}第一行：一个OFC神经元的活动，按当前形状的证据权重（即单个形状所携带的证据）的四分位数分组（颜色标示），活动对齐于形状开始呈现的时刻，并向后延伸至下一个形状呈现后100毫秒。该神经元的活动随当前形状的证据权重系统变化，编码单条证据的价值信息。第二行：一个DLPFC神经元的活动，按到当前时刻为止所有已呈现形状的累积证据权重的四分位数分组（红色深浅标示）。该神经元的活动随累积证据的增加而系统增强，表明其跨时间整合多条证据以形成判断。两图中虚线竖线标示形状在300毫秒处消失。改编自Lin et al. (2020)\cite{lin2020evidence}。}{\textbf{(A) Probabilistic reasoning task.} While the monkey fixated a central fixation point, four shapes were presented sequentially on the screen. In each trial, the four shapes were drawn randomly with replacement from a set of 10, each assigned a unique weight (see inset). Positive weights meant that the red target had a higher reward probability, whereas negative weights meant that the green target had a higher reward probability. The sum of the four shape weights determined the relative reward probabilities of the red and green targets. After the fixation point disappeared, the monkey made a saccade to either the red or the green target and received reward probabilistically. The red, yellow, green, and blue shadings on the time axis indicate the four sequential shape presentation periods.\\\textbf{(B) Example OFC and DLPFC neurons.} Top row: activity of an OFC neuron. Trials were divided into quartiles based on the evidence weight of the current shape (i.e., the evidence carried by a single shape), indicated by color. Response averages were aligned to shape onset and extended 100 ms into the next shape period. The neuron's activity varied systematically with the evidence weight of the current shape, encoding the value of individual pieces of evidence. Bottom row: activity of a DLPFC neuron. Trials were divided into quartiles by the cumulative evidence weight of all shapes presented up to that moment, indicated by the redness of the trace. The neuron's activity increased systematically as the cumulative evidence grew, indicating cross-temporal integration of multiple pieces of evidence to form a decision. Dashed vertical lines in both rows mark shape offset at 300 ms. Adapted from Lin et al. (2020)\cite{lin2020evidence}.}

DLPFC的第二项核心功能是认知控制——通过调控价值系统实现延迟满足与冲动抑制。针对DLPFC如何实施这一控制，已有研究提出了至少3种不同的机制。第一种机制是价值竞争。 McClure等人将跨期决策描述为双估值系统的竞争：多巴胺驱动的边缘回路（limbic circuit）偏好即时奖励，而以DLPFC为核心的额顶网络（frontoparietal network）则独立编码长远收益，决策结果取决于二者的相对激活强度\cite{mcclure2004separate}。在这一框架下，DLPFC通过直接输出与长远收益相关的价值信号，来对抗边缘系统所驱动的即时诱惑。第二种机制是独立于价值的控制。 Figner等人采用重复经颅磁刺激（repetitive transcranial magnetic stimulation, rTMS）短暂抑制左侧DLPFC，发现被试对即时选项的偏好显著增加，但对延迟奖励的主观估值并未发生改变\cite{figner2010lateral}。这表明DLPFC提供的是一种独立于估值过程的执行控制信号——它并不重新计算选项的价值，而是在价值比较完成之后施加自上而下的调控。第三种机制是调控价值表征以施加控制。Hare等人发现，vmPFC在所有被试中均编码目标价值，且该信号与自我控制程度无关；DLPFC的激活虽与价值评分无关，却与自我控制能力显著相关——DLPFC越活跃，vmPFC对“好吃但不健康”食物的响应就越低\cite{hare2009control}。这一结果提示，DLPFC虽不直接参与估值计算，但可以通过下调vmPFC中诱惑性选项的价值来间接实现控制。上述三种假说在具体机制上存在分歧——分别强调价值竞争、独立控制信号与估值权重调节——但共同指向同一结论：DLPFC是价值决策网络中的关键调控节点，其功能状态能够系统地影响价值系统的最终输出。

\subsection{ACC：行动执行前的广泛预备性整合}\label{chap1:ACC}
ACC的背侧部（dorsal anterior cingulate cortex, dACC）被视为将动机与认知状态转化为行动指令的关键枢纽，其核心区域为布罗德曼24区\cite{paus2001primate}（Brodmann area 24, BA 24）。dACC极少接受初级感觉皮层（primary sensory cortices）与颞叶联合皮层（temporal association cortex）的直接投射，与记忆相关的海马旁区域也几乎不发生联系\cite{heilbronner2016dorsal}；来自边缘结构（limbic structures）的情感信号需经腹侧ACC（ventral anterior cingulate cortex, vACC）间接中转，方能抵达dACC的运动相关区域\cite{paus2001primate}。与匮乏的感觉输入形成鲜明对照的是，dACC与运动相关脑区有广泛连接：与初级运动皮层（primary motor cortex, M1）、前运动皮层及辅助运动区（supplementary motor area, SMA）保持紧密的双向连接；其尾侧区域（扣带运动区，cingulate motor areas, CMAs）还直接发出皮质脊髓束（corticospinal tract）参与运动执行\cite{devinsky1995contributions, paus2001primate}。这种“弱感觉、强运动”的连接模式，使得dACC可以被理解为“前前运动皮层”（pre-premotor cortex）——它在行动执行之前负责广泛的预备性整合，涵盖情境评估与策略调控等功能\cite{heilbronner2016dorsal}。整体而言，dACC所编码的信息虽最终指向运动输出，但其抽象程度远高于运动执行本身。

ACC的功能之一是参与价值决策，尤其在需要考虑与行动相关的成本时更为突出。Balewski等人发现，在决策形成过程中，ACC的神经元呈现类似证据累积的爬升活动（ramping activity），其爬升速率受OFC所表征的价值状态调节——当OFC表征的选项价值更高时，ACC的爬升更快\cite{balewski2023value}。这表明ACC可能持续累积来自OFC的价值信号以作出决策。Cai等人的研究发现，dACC神经元在经济选择中编码被选价值（chosen value）和被选类型（chosen juice），同时整合运动方向信息，提示dACC处于将价值转化为行动的关键节点（图~\ref{fig0_5:cai2012}）\cite{cai2012neuronal}。值得注意的是，该研究未观察到ACC编码选项价值（offer value），一个可能的原因是实验任务未引入显著的体力成本。

\begin{figure}[H]
    \centering
    \includegraphics[width=\textwidth]{Figure0_5_Cai2012}
    \bicaption{\enspace dACC神经元编码主观价值与选择信号}{\enspace dACC neurons encode subjective value and choice signals}
    \label{fig0_5:cai2012}
\end{figure}
\fignotetext{本实验采用与图~\ref{fig0_3:padoa2006}A相同的实验范式。横轴为选项类型（按B:A比值排列），黑色散点及S形拟合曲线为该神经元记录当天的行为选择模式，左上角标注了拟合得出的相对价值（如1A=3.2B）；彩色散点为神经元放电率，红/绿色分别表示右向/左向眼跳，菱形/圆形分别表示选择了果汁A/果汁B。左、中、右三列分别为编码被选价值、动作域选择（左--右）、物品域选择（果汁A--果汁B）的神经元。改编自Cai \& Padoa-Schioppa (2012)\cite{cai2012neuronal}。}{The same task paradigm as in Fig.~\ref{fig0_3:padoa2006}A was used. The x-axis represents offer type ranked by the B:A ratio. Black dots and the sigmoid fit show the behavioral choice pattern from the same recording day, with the relative value from the fit indicated at the top left (e.g., 1A=3.2B). Colored dots represent neuronal firing rates; red and green indicate rightward and leftward saccades, respectively, while diamonds and circles indicate trials in which the monkey chose juice A and juice B, respectively. The left, middle, and right columns show neurons encoding chosen value, motor-domain choice (left vs. right), and goods-domain choice (juice A vs. juice B), respectively. Adapted from Cai \& Padoa-Schioppa (2012)\cite{cai2012neuronal}.}

当任务包含真实的体力成本时，ACC的价值编码明显增强：Kennerley等人在实验中同时操纵奖励量（magnitude）、奖励概率（probability）与体力成本（effort）三种决策变量，发现ACC对价值信号的选择性响应在编码强度与广度上均超过OFC和DLPFC\cite{kennerley2009neurons}。Prévost等人比较了延迟与体力两类成本的神经表征，发现vmPFC与腹侧纹状体（ventral striatum，VS）负责处理延迟折扣，而ACC与前脑岛（anterior insula）构成的网络则专门处理体力成本相关的估值\cite{prevost2010separate}。进一步的电刺激实验证实了ACC在运动成本相关价值决策中的因果贡献\cite{amemori2012localized}。综合上述证据，ACC在涉及体力成本的价值决策中居于关键地位。但也需指出，相关证据并非完全一致。例如，Hosokawa等人在猕猴电生理实验中也比较了延迟与体力两类成本，发现ACC是唯一在两种成本场景下均显著编码价值的脑区\cite{hosokawa2013mechanisms}，这与Prévost等人在人类被试中观察到的延迟（OFC/vmPFC）与体力（ACC-insula）之间的分离结果相矛盾。两类结果之间的分歧提示，ACC在价值决策中的具体功能可能受任务设计、物种差异或分析方法的影响，其确切作用仍有待进一步检验。

ACC的另一项重要功能是在觅食与探索中编码替代选项的价值，从而指导“停留-切换”（stay-switch）决策。Hayden等人设计了一个虚拟觅食任务：猕猴可选择停留在当前环境中持续收获（收益递减），或切换至新环境（收益重置但需付出时间成本）。猕猴的行为符合边际价值定理（Marginal Value Theorem）——当停留在当前环境的收益比切换到新环境所能获得的预期收益低时，动物即付出时间成本迁往新环境。实验中dACC神经元的发放率随猕猴在当前环境中的停留时间持续升高，累积至阈值时触发切换行为；当切换的时间成本提高时，该阈值亦升高\cite{hayden2011neuronal}。这表明dACC编码的是替代选项的相对价值（即离开的动机强度），而非当前即时收益，并据此驱动“停留-切换”决策。Kolling等人的人类fMRI研究观察到相似的结果：在搜索阶段，需要决定是否放弃当前选项去探索替代选项时，ACC的活动与环境中替代选项的平均价值正相关，而与当前选项的价值无关；与之形成对照，vmPFC在搜索阶段不编码替代选项的价值，仅在随后的价值决策阶段编码选项间的价值差异\cite{kolling2012neural}。Tervo等人通过光遗传学手段揭示了更为复杂的因果结构：在大鼠评估当前策略是否值得坚持的“重评估点”（re-evaluation point）抑制ACC，大鼠更倾向于切换至替代选项；而在大鼠遭遇具体替代选项、即将做出接受决定的“承诺点”（commitment point）抑制ACC，大鼠反而更难切换\cite{tervo2021anterior}。这意味着ACC并非简单地促进或抑制切换行为，而是在决策过程的不同阶段承担着不同的计算功能。

需要指出的是，主观价值的神经表征并不局限于以上脑区。腹侧纹状体\cite{kable2007neural, lebreton2009automatic, demartino2009neurobiology, peters2009overlapping, prevost2010separate, bartra2013valuation}和后扣带回（posterior cingulate cortex, PCC，Kable等人、Lebreton等人、Peters等人和Grueschow等人\cite{kable2007neural, lebreton2009automatic, peters2009overlapping, grueschow2015automatic}）等区域同样被发现编码主观价值信号。综合来看，价值编码并非依赖单一的“价值计算器”，而是通过一个分布式网络实现。

\subsection{财富的神经表征及其对价值编码的调控}\label{chap1:wealth_neural}
多项研究表明，财富水平的神经表征分布于ACC、vmPFC、前岛叶皮层（anterior insular cortex, AIC）等多个脑区。Seo等人在猕猴博弈任务中发现，背内侧额叶皮层（dorsal medial frontal cortex, DMFC）与dACC中显著追踪当前持有资产（代币）的神经元比例（分别为64.5\%与57.3\%）远高于DLPFC（23.7\%），提示包括dACC在内的内侧额叶在基于代币的行为调整中可能起着关键作用、\cite{seo2009behavioral}。Juechems等人在连续风险决策任务中发现，随着财富升高，人类被试的风险偏好系统性下降；在此过程中，vmPFC能够在不依赖外部信号的情况下隐式追踪累积奖励\cite{juechems2017ventromedial}。Yang等人在猕猴代币赌博任务中发现了AIC的双重作用：首先，AIC神经元以参数型（parametric，单调增或单调减，连续变化）、类别型（categorical，高/低两类）与数值型（numerical，偏好某一位置，向两侧递减）等不同方式编码当前财富水平；同时，对于同一选项（即相同的奖赏数额），AIC神经元编码的主观价值强度随财富水平升高而系统性降低，与效用理论中边际价值递减的预测一致（图~\ref{fig0_6:yang2022neural}）\cite{yang2022primate}。

Nguyen等人对前额叶六个区域的系统比较揭示了更精细的功能分工：前扣带皮层腹侧部（ventral anterior cingulate cortex, vACC）在群体层面选择性地编码参考点本身（即当前累积财富）；而dACC与DLPFC则分别编码预期价值与已实现结果的主观价值，且两者的调谐曲线斜率随财富增加而下降，即对价值的编码增益随财富水平升高而系统性衰减，符合效用理论中边际价值递减的预测\cite{nguyen2025neural}。

\begin{figure}[H]
    \centering
    \includegraphics[width=\textwidth]{Figure0_6_Yang2022Neural}
    \bicaption{\enspace AIC神经元编码代币财富与财富依赖的主观价值}{\enspace AIC neurons encode token wealth and wealth-dependent subjective value}
    \label{fig0_6:yang2022neural}
\end{figure}
\fignotetext{\textbf{(A) AIC神经元对起始代币数的三种编码模式。}实验任务同图~\ref{fig0_2:yang2022}所示的代币赌博任务。左：参数型编码，神经元放电率随起始代币数单调变化（递增或递减），以连续方式编码财富水平。中：类别型编码，放电率将起始代币数划分为高、低两类，两类之间反应显著不同。右：数值型编码，神经元对特定代币数值（如此例中的4）有最大反应，向两侧递减，呈调谐曲线形态。\\\textbf{(B) AIC神经元对选项期望价值的编码强度随财富水平的变化。}分析基于增益价值神经元（仅在获益情境编码价值，$n=12$）和损失价值神经元（仅在损失情境编码价值，$n=39$）两类群体。对每个神经元，首先将其在强制选择试次按起始代币数分为低财富组（0--2）和高财富组（3--5），在每组内以选项期望价值（结果$\times$概率）为自变量、神经元放电率为因变量做线性回归，得到标准化回归系数（SRC），用于衡量该神经元在该财富条件下对价值变化的敏感度。随后，分别计算增益价值神经元和损失价值神经元在低/高财富条件下SRC的均值，得到四条平均斜率，据此绘制图中四条回归线：绿色为增益价值神经元，红色为损失价值神经元；浅色为低财富条件，深色为高财富条件；横轴为选项期望价值（EV），纵轴为归一化放电率；阴影区域为$\pm$标准误。无论是增益情境还是损失情境，高财富条件下的回归线斜率均比低财富条件下更平缓，表明当持有更多代币时，AIC神经元对同一期望价值的编码强度系统性降低，符合边际效用随财富积累递减的预期。改编自Yang et al. (2022)\cite{yang2022primate}。}{\textbf{(A) Three encoding modes of start token number by AIC neurons.} The task was the same token-based gambling task as in Fig.~\ref{fig0_2:yang2022}. Left: parametric encoding, in which the neuron's firing rate varied monotonically (increasing or decreasing) with the start token number, encoding wealth level in a continuous manner. Middle: categorical encoding, in which the firing rate distinguished between low and high start token levels, with significantly different responses between the two categories. Right: numerical encoding, in which the neuron was maximally responsive to a specific token number (e.g., 4 in this example) and tuned away from this preferred value toward both smaller and larger numbers.\\\textbf{(B) Changes in AIC neuronal sensitivity to expected option value as a function of token wealth.} The analysis was based on two independently identified neuronal populations: Gain-value neurons (encoding value only in the gain context, $n = 12$) and Loss-value neurons (encoding value only in the loss context, $n = 39$). For each neuron, forced-choice trials were first split into low start-token trials (0--2 tokens) and high start-token trials (3--5 tokens). Within each subset, a linear regression was performed with the expected value of the option (outcome $\times$ probability) as the predictor and the neuron's firing rate as the response. The resulting slope was standardized (standardized regression coefficient, SRC) to quantify the neuron's sensitivity to value change under that wealth condition. The mean SRCs were then computed separately for Gain-value and Loss-value neurons at low and high token levels, yielding four average slopes. Based on these, four regression lines were plotted: green for Gain-value neurons, red for Loss-value neurons; light colors for low-wealth conditions, dark colors for high-wealth conditions. The x-axis represents expected option value (EV), and the y-axis represents normalized firing rate; shaded regions indicate $\pm$ SEM. In both gain and loss contexts, the regression slopes were shallower under high-wealth conditions than under low-wealth conditions, indicating that as token wealth increased, AIC neurons' sensitivity to expected value decreased systematically, consistent with the diminishing marginal utility of wealth. Adapted from Yang et al. (2022)\cite{yang2022primate}.}

Ferro等人的研究发现随着代币数量接近奖励兑换阈值，猕猴的风险偏好系统性降低；同时dACC中相当大比例的神经元持续追踪当前代币数量，据此动态调整参考点，以参考依赖的方式重新标定选项的主观价值\cite{ferro2026accumulation}。

综合来看，财富的神经表征与主观价值的编码共享“分布式”这一基本特征：财富并非游离于价值系统之外的独立变量，而是作为动态参照点深度嵌入这一网络，通过调节各节点的编码增益参与经济决策。


\section{研究概览}\label{chap1:overview}
然而，现有代币范式存在两方面局限。第一，固定的代币兑换阈值（如普遍在累积到6枚代币之后兑换奖励）在任务中引入了外生财富参考点。若要在该类范式中准确解析财富效应，需额外建模以剥离阈值影响\cite{ferro2026accumulation}；若研究旨在系统考察参考点本身的作用，则应采用专门的实验设计（如将兑换阈值设置为实验变量），而非将其作为混淆因素遗留于任务结构之中。第二，现有实验中的选项始终在代币得失这一维度上进行比较，考察的是同一价值类型内部的权衡，无法体现财富作为通用媒介在不同类型价值之间调配资源的基本属性。

赚取–消费范式可同时回应上述两个问题。在每个试次中，动物需在两种选项间做出选择：赚取代币，或消费已累积的代币以换取即时奖励。这一选择发生于次级强化物（代币）与初级强化物（饮用水）之间，构成跨类别的价值权衡。这种“赚取–消费”结构也正是经济生活中最普遍的决策形式——例如加班与请假的权衡，本质上就是薪资与休闲之间的竞争。人类行为研究已显示，财富水平对此类决策具有系统性影响，但在非人灵长类中，类似行为鲜有研究。

在理论建模层面，此前的代币任务中选项天然具有得失结构——赚取代币为收益，失去代币为损失——因此常以前景理论建模，以同时处理收益域与损失域的价值函数。赚取–消费决策则不同：赚取代币是资产积累，消费代币换取奖励同样构成收益，两个选项均不涉及损失域。因此，本研究抛弃复杂的参考点设定与损失域建模，转而采用期望效用理论，核心假设为：赚取选项的价值随着财富积累下降，消费选项的相对吸引力随之上升。

基于上述考量，本研究训练三只恒河猴完成赚取–消费抉择任务。行为结果表明，动物的选择偏好受到选项价值与财富水平的共同调控——随着财富积累，消费倾向显著上升，赚取倾向显著下降，符合边际效用递减的定量预测。据我们所知，这是首次在非人灵长类中研究财富依赖的赚取–消费决策。同步记录的前额叶三个子区（ACC、OFC、DLPFC）呈现出差异化的编码模式：ACC对财富状态及其动态更新的表征最为突出；OFC在编码财富的同时，也编码经财富调制的主观价值与最终选择信号；DLPFC则主要编码动作层面的选择信息。这些结果表明，灵长类前额叶皮层通过子区域间的功能分工支持财富依赖的经济决策，为理解财富如何塑造经济决策行为提供了单神经元层面的证据。